
\subsubsection{System Design}

The AI documentation copilot operates in three stages: (1) analyze uploaded dataset properties to extract contextual features, (2) retrieve relevant examples from a curated corpus of high-quality documentation, and (3) generate contextual suggestions via large language model few-shot prompting.

\textbf{Dataset Property Analysis:} The system extracts file formats, data types, size metrics, and statistical distributions. Domain classification identifies dataset type (vision, NLP, or tabular) based on file extensions, column names, and content patterns.

\textbf{Few-Shot Prompting:} The system uses Llama-3-8B-Instruct with few-shot prompting rather than fine-tuning. For each documentation field, the system retrieves 3-5 relevant examples from a corpus stratified by domain (vision, NLP, tabular) and constructs prompts including dataset properties, the documentation section to fill, and retrieved examples. Generation uses temperature 0.7 and maximum length 500 tokens with GPU acceleration.

\textbf{Implementation Note:} This proof-of-concept validates the few-shot prompting mechanism and user interaction design. The pilot used a representative corpus structure to validate the approach. Production deployment would require systematic curation and validation of the complete example corpus as described in the architecture.

\textbf{User Interaction:} Suggestions appear inline as researchers fill documentation template fields. Each suggestion presents three action buttons: Accept (insert as-is), Modify (accept into editable field for refinement), and Reject (dismiss and write from scratch). All user actions are logged with timestamps.

\subsubsection{Deployment Protocol}

The pilot ran for 2 weeks with 75 researchers recruited as volunteers on HuggingFace. Each participant documented 1-5 datasets during the period, generating approximately 25 suggestions per user (1,875 total suggestions). All suggestion interactions were logged to JSON files for analysis.

\textbf{Participant Recruitment:} Volunteers who opted into the pilot study. Self-selection likely inflates acceptance rates compared to the general population but provides initial feasibility evidence.

\textbf{Evaluation Metric:} Median per-user acceptance rate, calculated as (accepted + modified) / total_suggestions × 100\%. Aggregation at user level (not suggestion level) avoids biasing toward high-volume users.

\textbf{Success Threshold:} We pre-registered 70\% median acceptance as the deployment feasibility threshold. Acceptance below 40\% would indicate mechanism failure; between 40-70\% would suggest promising but immature technology; above 70\% validates deployment readiness for the engagement mechanism.

\textbf{Stratified Analysis:} Acceptance rates were analyzed by dataset type (vision, NLP, tabular) to evaluate cross-domain generalization.
